\begin{table}[htbp]
\caption{Prevalencia ponderada de depresión por piso altitudinal Estimación imprecisa (CV > 0,20): la muestra de esta categoría es pequeña. Debe leerse por su intervalo de confianza, no por la cifra puntual. Afecta a 1 de 5 filas, marcadas con asterisco.}
\label{tab:tabla_33_prevalencia_por_altitud}
\begin{tabular}{lrrrrrrrrl}
\toprule
altitud\_categoria & prevalencia & se & ic\_inf & ic\_sup & cv & deff & df & n & precision \\
\midrule
Costa / selva baja (0-499 m) & 6.673000 & 0.351000 & 6.017000 & 7.394000 & 0.053000 & 2.885000 & 1568 & 14615 &  \\
Yunga (500-1 999 m) & 5.955000 & 0.549000 & 4.962000 & 7.131000 & 0.092000 & 2.541000 & 392 & 4715 &  \\
Quechua (2 000-2 999 m) & 7.109000 & 0.559000 & 6.086000 & 8.290000 & 0.079000 & 1.705000 & 313 & 3610 &  \\
Suni (3 000-3 999 m) & 8.898000 & 0.552000 & 7.872000 & 10.043000 & 0.062000 & 2.454000 & 533 & 6538 &  \\
Puna (4 000 m o más) & 11.702000 & 2.598000 & 7.419000 & 17.978000 & 0.222000 & 4.879000 & 59 & 748 & * imprecisa (CV > 0,20) \\
\bottomrule
\end{tabular}
\end{table}
